\section{自动微分：数值与梯度的分别实现}
\label{est:sec:kl-autodiff}
前面的统一框架估计已经指定的积分。本节检查标量程序的反传究竟对应哪个积分。若一个程序同时用于 KL 日志与训练，审计对象应记录为 forward 数值和 backward 向量组成的有序对。
本节涉及随机梯度的期望与偏差时默认 $\beta=P$；固定索引的逐样本求导规则本身不依赖抽样分布。

\subsection{一个函数的两种求导路径}

\begin{theorem}[固定样本反传与采样测度补偿]
\label{est:thm:kl-autodiff-transform}
令 $\phi$ 是不另依赖参数 $\theta$ 的可微实函数，$A\sim P$，反传时固定样本索引。定义
\begin{equation*}
 F_A=\phi(\ell_A),\qquad
 c_A=\frac{P(A)}{\sg(P(A))},\qquad
 L_A=c_A\phi(\ell_A).
\end{equation*}
两者的 forward 值逐样本相同，固定样本导数分别为
\begin{equation*}
 \nabla F_A=\phi'(\ell_A)\psi_A,\qquad
 \nabla L_A=\{\phi(\ell_A)+\phi'(\ell_A)\}\psi_A.
\end{equation*}
后者的期望等于 $\nabla I_P(\phi(\ell))$。若仅将系数停止梯度，$c_A\sg(\phi(\ell_A))$ 的导数则为 $\phi(\ell_A)\psi_A$。
\end{theorem}

\begin{proof}
固定索引下 $\nabla \ell_A=\psi_A$，链式法则给出第一式。forward 时 $c_A=1$；分母停止梯度，因而 $\nabla c_A=\psi_A$。乘积法则给出第二式。再由式~\eqref{est:eq:kl-measure-derivative}，
\begin{equation*}
 \nabla I_P(\phi(\ell))
 =I_P\bigl(\phi'(\ell)\psi+\phi(\ell)\psi\bigr),
\end{equation*}
等于第二种导数的期望。停止整个 $\phi$ 时只保留 $\nabla c_A$ 一项，得到最后的规则。
\end{proof}

这个变换把 K1--K4 的差异归到同一个问题：给定 forward 函数后，程序保留了哪些乘积法则项。数值无偏是 $I_P(\phi(\ell))=K$，裸反传无偏则是 $I_P(\phi'(\ell)\psi)=g$，两者是不同等式。

\begin{center}
\begin{tabular}{lll}
\toprule
程序 & forward 值 & 固定样本 backward\\
\midrule
K1 & $\ell_A$ & $\psi_A$\\
K2 & $\ell_A^2/2$ & $\ell_A\psi_A$\\
K3 & $e^{-\ell_A}+\ell_A-1$ & $(1-e^{-\ell_A})\psi_A$\\
K4 & $\ell_A$ & $\ell_A\psi_A$\\
K4$^{\mathrm{live}}$ & $\ell_A$ & $(\ell_A+1)\psi_A$\\
K3$^{++}$ & $e^{-\ell_A}+\ell_A-1$ & $\ell_A\psi_A$\\
\bottomrule
\end{tabular}
\end{center}
数值方差按第二列计算，梯度方差按第三列计算。K5 的目标是$\KL(Q\Vert P)$，随后单独代入同一求导规则。

\subsection{K4 与固定样本的梯度替换}
K1 与 K3 的数值都无偏；K2 是以二阶矩换取局部稳定性的有偏 proxy。若需要同时让数值和自动微分梯度对应同一个 $K$，使用停止 log-ratio 系数的 K4：
\begin{equation*}
 c_A=\frac{P(A)}{\operatorname{sg}(P(A))},\qquad
 \widehat K_4=c_A\sg(\ell_A).
\end{equation*}
数值上 $c_A=1$，所以 $\E[\widehat K_4]=K$；反传时
\begin{equation*}
 \nabla_\theta\widehat K_4=\ell_A\psi_A,
\end{equation*}
因此
\begin{equation*}
 \E_P[\nabla_\theta\widehat K_4]=\nabla_\theta K=g.
\end{equation*}
K4 与 K1 的 forward 值逐样本相同，故 K4 的数值偏差为零，数值方差也为 $\Var_P(\ell)$。stop-gradient 只改变求导规则。另定义本项目单 token 测试采用的 live-log-ratio 变体
\begin{equation*}
 \widehat K_4^{\mathrm{live}}=c_A \ell_A,\qquad
 \nabla\widehat K_4^{\mathrm{live}}=(\ell_A+1)\psi_A.
\end{equation*}
它的数值仍与 K4 相同，梯度期望也相同，但梯度方差不同。后文始终保留上标 $\mathrm{live}$，不将两种反传路径共用一个名称。

\begin{theorem}[K4 的数值与梯度无偏性]
\label{est:thm:k4}
在本章正支持与固定参考假设下，$\widehat K_4$ 和 $\widehat K_4^{\mathrm{live}}$ 的数值均对 $K$ 无偏，自动微分梯度均对 $g$ 无偏。
\end{theorem}
\begin{proof}
两者的数值为 $\ell_A$，因此
\begin{equation}
\E_P[\widehat K_4]=\E_P[\widehat K_4^{\mathrm{live}}]=\sum_vP(v)\ell_v=K.
\end{equation}
两者的梯度分别为 $\ell_A\psi_A$ 与 $(\ell_A+1)\psi_A$。它们的期望之差为
\begin{equation}
\sum_vP(v)[(\ell_v+1)\psi_v-\ell_v\psi_v]=\sum_vP(v)\psi_v=0.
\end{equation}
前者的期望是 $I_P(\ell\psi)=g$，所以后者的期望同样等于 $g$。
\end{proof}

\subsection{K3$^{++}$ 与$\KL(Q\Vert P)$ 的 K5}

EMA-PG 的公开实现使用上述停止系数 K4，并定义 K3$^{++}$ 与 K5；以下公式先去掉 clipping，单独分析数学对象。\footnote{代码版本固定为 EMA-PG \path{f9b9776}，定义见该版本的 \href{https://github.com/LunjunZhang/ema-pg/blob/f9b977632ad948ce935432a517c22291bfaba562/search/verl/trainer/ppo/core_algos.py}{\path{core_algos.py}}。}
\begin{equation*}
 \widehat K_{3^{++}}=c_A(e^{-\ell_A}+\ell_A-1).
\end{equation*}
它的数值与 K3 相同，因此对$K=\KL(P\Vert Q)$ 无偏，数值方差也与 K3 相同。其反传为
\begin{equation*}
 \nabla\widehat K_{3^{++}}
 =\{e^{-\ell_A}+\ell_A-1+1-e^{-\ell_A}\}\psi_A=\ell_A\psi_A.
\end{equation*}
因此它与 K4 的 backward 逐样本相同，梯度偏差为零，梯度方差也完全相同；数值方差的差异不能解释成梯度方差的差异。

K5 对应$\KL(Q\Vert P)$ 目标 $K^{\mathrm F}=\KL(Q\Vert P)$，定义为
\begin{equation*}
 \widehat K_5=-\sg(e^{-\ell_A})\ell_A+\log P(A)-\sg(\log P(A)).
\end{equation*}
其 forward 值为 $-e^{-\ell_A}\ell_A$。在 $A\sim P$ 下，
\begin{equation*}
 \E_P[\widehat K_5]=-\sum_vQ(v)\ell_v=K^{\mathrm F},\qquad
 \Var(\widehat K_5)=\E_P[e^{-2\ell_A}\ell_A^2]-(K^{\mathrm F})^2.
\end{equation*}
它相对本章默认目标 $K=\KL(P\Vert Q)$ 的数值偏差是 $K^{\mathrm F}-K$，不能与 K1--K4 按同一个无偏目标直接并列。反传为
\begin{equation*}
 \nabla\widehat K_5=(1-e^{-\ell_A})\psi_A,\qquad
 \E_P[\nabla\widehat K_5]=-\sum_vQ(v)\psi_v=\nabla K^{\mathrm F}.
\end{equation*}
其梯度方差为
\begin{equation*}
 \E_P[(1-e^{-\ell_A})^2\lVert \psi_A\rVert^2]
 -\left\lVert\sum_vQ(v)\psi_v\right\rVert^2.
\end{equation*}
这里$\KL(Q\Vert P)$ 目标只改变 KL 的积分方向，student 与 score 的定义保持不变。

若数值稳定化把 K3 的函数 $\phi(\ell)=e^{-\ell}+\ell-1$ 改成分段可导的 $\widetilde\phi(\ell)$，K3$^{++}$ 的数值偏差和梯度偏差分别变成
\begin{equation*}
 \E_P[\widetilde\phi(\ell_A)]-K,\qquad
 \E_P[(\widetilde\phi(\ell_A)+\widetilde\phi'(\ell_A)-\ell_A)\psi_A].
\end{equation*}
公开代码对 log-ratio 与输出均有 clamp，因此无 clipping 恒等式的适用条件是相关 clamp 未激活，或相关偏差项经额外核验为零。

\subsection{数值 estimator 的自动微分梯度}

把已经抽到的动作视为固定索引，对 K1--K4 直接反传，得到
\begin{equation*}
 G_1=\psi_A,\quad G_2=\ell_A\psi_A,\quad
 G_3=(1-e^{-\ell_A})\psi_A,\quad G_4=\ell_A\psi_A.
\end{equation*}
令 $u_Q=\E_Q[\psi_A]$，则这些随机梯度相对于 $g$ 的偏差为
\begin{equation*}
 \operatorname{Bias}(G_1)=-g,\quad
 \operatorname{Bias}(G_2)=0,\quad
 \operatorname{Bias}(G_3)=-u_Q-g,\quad
 \operatorname{Bias}(G_4)=0.
\end{equation*}
对应方差逐项为
\begin{align*}
 \Var(G_1)&=\E_P\lVert \psi_A\rVert^2,\\
 \Var(G_2)&=\E_P[\ell_A^2\lVert \psi_A\rVert^2]-\lVert g\rVert^2,\\
 \Var(G_3)&=\E_P[(1-e^{-\ell_A})^2\lVert \psi_A\rVert^2]-\lVert u_Q\rVert^2,\\
 \Var(G_4)&=\E_P[\ell_A^2\lVert \psi_A\rVert^2]-\lVert g\rVert^2.
\end{align*}
例如 $\E_P[G_3]=\E_P[\psi_A]-\E_P[(Q(A)/P(A))\psi_A]=-u_Q$。由此得到两个明确区别：K2 的数值有偏但其固定样本反传梯度无偏；K1 与 K3 的数值无偏但其裸反传梯度通常有偏。这些期望在每次抽样均使用当前 $P$ 时成立；不能将固定样本的导数与包含采样测度导数的 $\nabla_\theta\E_{P_\theta}$ 混用。

确定性 Top-$J$ 数值目标的梯度也可以显式计算。记
\begin{equation*}
 d_{P_S}=\nabla_\theta P_S=\sum_{v\in S}P(v)\psi_v,
 \qquad \ell_{\mathrm{mass}}=\log(\tau_P/\tau_Q).
\end{equation*}
由乘积法则与 $\nabla_\theta\tau_P=-d_{P_S}$，有
\begin{align*}
 \nabla_\theta H&=h_g+d_{P_S},\\
 \nabla_\theta(H+M)&=h_g-\ell_{\mathrm{mass}}d_{P_S}=h_g+\tau_Pb_{\mathrm{TA}},\\
 \nabla_\theta K_J^{\mathrm{norm}}
 &=\frac{h_g}{P_S}-\frac{H}{P^2}d_{P_S}.
\end{align*}
所以三者相对于完整 KL 梯度的偏差依次为
\begin{equation*}
 d_{P_S}-\tau_P\mu_g,\qquad
 \tau_P(b_{\mathrm{TA}}-\mu_g),\qquad
 \frac{h_g}{P_S}-\frac{H}{P^2}d_{P_S}-g,
\end{equation*}
方差均为零。尤其 $h_g$ 是停止系数的头部 score 梯度，$\nabla H=h_g+d_{P_S}$ 是头部数值的真实梯度；两者必须分别命名。

\subsection{停止梯度 surrogate 的显式构造}

为把数学上的梯度向量对应到可反传的标量，记 $\operatorname{sg}$ 为停止梯度算子。固定抽样动作和事件，定义
\begin{align*}
 L_{\mathrm{HETS}}
 &=\sum_{v\in S}\operatorname{sg}(P(v)\ell_v)\log P(v)
   +J_A\operatorname{sg}(\ell_A)\log P(A),\\
 L_{\mathrm{Res}}
 &=L_{\mathrm{HETS}}
   +\operatorname{sg}((\tau_P-J_A)\ell_{\mathrm{mass}})\log\tau_P.
\end{align*}
逐项反传，得到
\begin{align*}
 \nabla L_{\mathrm{HETS}}&=h_g+J_A\ell_A\psi_A,\\
 \nabla L_{\mathrm{Res}}
 &=h_g+J_A\ell_A\psi_A+(\tau_P-J_A)\ell_{\mathrm{mass}}\nabla\log\tau_P\\
 &=h_g+\tau_Pb_{\mathrm{TA}}+J_A(\ell_A\psi_A-b_{\mathrm{TA}}).
\end{align*}
这两条等式依赖停止整个系数的梯度，包括其中的 $P(v)$、$\ell_v$、$\tau_P$ 和 $\ell_{\mathrm{mass}}$；若遗漏停止操作，就会引入额外导数。

若希望 forward 展示某个数值估计量 $F$，同时 backward 使用 $L$ 的梯度，可定义
\begin{equation*}
 \widetilde L=\operatorname{sg}(F)+L-\operatorname{sg}(L).
\end{equation*}
逐样本有 $\widetilde L=F$，但 $\nabla\widetilde L=\nabla L$。因此 forward 的数值偏差和 backward 的梯度偏差可以独立设计，审计时必须同时给出 $F$ 和 $L$。

\subsection{K4 尾部反传与 HETS 的精确关系}
\label{est:sec:kl-k4-hets}

前文 HETS 数值估计量在项目中也称为 EMA，其公式为 $\widehat K^{\mathrm{HETS}}=H+J_A\ell_A$。若头部按真实数值 $H$ 求导，而尾部使用本章明确命名的 K4$^{\mathrm{live}}$ 构造，则自动微分得到
\begin{equation*}
 \widehat g^{\mathrm{EMA\text{-}K4}}
 =h_g+d_{P_S}+J_A(\ell_A+1)\psi_A,\qquad d_{P_S}=\nabla_\theta P_S.
\end{equation*}
这与 HETS 的 $h_g+J_A\ell_A\psi_A$ 相差 $d_{P_S}+J_A\psi_A$。因为 $\E_P[J_A\psi_A]=\nabla_\theta\tau_P=-d_{P_S}$，差值期望为零，但样本和方差均不同。

\begin{theorem}[头尾 score 常数的匹配]
对固定实数 $\upsilon$，定义
\begin{equation*}
 \widehat g_\upsilon
 =h_g+\upsilon d_{P_S}+J_A(\ell_A+\upsilon)\psi_A.
\end{equation*}
则它对 $g$ 无偏，且
\begin{equation}
 \Var(\widehat g_\upsilon)
 =\sum_{v\notin S}P(v)(\ell_v+\upsilon)^2\lVert \psi_v\rVert^2
  -\lVert\tau_P\mu_g-\upsilon d_{P_S}\rVert^2.
 \label{est:eq:kl-alpha-var}
\end{equation}
$\upsilon=0$ 是 HETS，$\upsilon=1$ 是真实头部数值加 masked K4$^{\mathrm{live}}$ 尾部的反传。
\end{theorem}

\begin{proof}
期望为
\begin{equation*}
 h_g+\upsilon d_{P_S}+\tau_P\mu_g+\upsilon\sum_{v\notin S}P(v)\psi_v
 =h_g+\tau_P\mu_g=g.
\end{equation*}
确定性头部不贡献协方差；随机尾部的均值为 $\tau_P\mu_g-\upsilon d_{P_S}$，二阶矩为式~\eqref{est:eq:kl-alpha-var} 的求和项，两者相减得到方差。
\end{proof}

如果误把真实头部梯度 $h_g+d_{P_S}$ 与尾部的 $J_A\ell_A\psi_A$ 拼接，得到的期望为 $g+d_{P_S}$，偏差为 $d_{P_S}$。若头部停止系数为 $h_g$，尾部却使用 $J_A(\ell_A+1)\psi_A$，偏差为 $-d_{P_S}$。两端的 score 常数必须匹配；仅凭“精确头部加 sampled tail”这句话无法判断无偏性。

EMA-PG 固定版本的 \path{compute_memory_efficient_kl} 在 \path{full_reverse} 路径中直接反传头部 $H$，尾部使用停止 log-ratio 的 K4。因此在共同固定支持、无 clipping 且尾部测度校正正确时，forward 无偏，backward 的偏差为 $d_{P_S}$。actor 传入的 log-importance weight 已 detach，不能通过它的导数消掉 $d_{P_S}$。\footnote{同一固定版本的 actor 路径见 \href{https://github.com/LunjunZhang/ema-pg/blob/f9b977632ad948ce935432a517c22291bfaba562/search/verl/workers/actor/dp_actor.py}{\path{dp_actor.py}}。}本章所证明的无偏 EMA 数值结构与论文 HETS 的向量结构，不能据名称相近就认定为该路径的逐样本反传行为。

本仓库单 token K4$^{\mathrm{live}}$ 的数值与梯度测试位于 \path{tests/test_kl_estimators.py}，对应 $(\ell_A+1)\psi_A$。

结构统计程序位于 \path{experiments/analysis/ema_topk_k_analysis.py}。它采用 $h_g+d_{P_S}$ 头部与 $\ell_A\psi_A$ 尾部，其期望具有上述 $d_{P_S}$ 偏差；这些统计不能作为正确匹配头尾常数的估计量或 HETS 的无偏性证据。

生产损失位于 \path{src/mirroropd/training/losses.py}。其中 OPD baseline 使用停止梯度的 Tail-aware 标量，属于 scalar baseline 路径。本章的数学分类分别描述这三个实际对象。
